\section{Colliding points and degenerating sections}
\label{sec:limits}

We turn the multiplicities at the remaining $q$ points into derivative
conditions at one point of $\Tot(A)$. Universal existence will force
these conditions to have a nonzero solution in $\W_\tau$. We then
choose bases whose derivative matrices have an explicit limit as
$\tau\to0^+$. Section~\ref{sec:interpolation} will prove full column
rank when the common multiplicity is sufficiently large at a fixed
ratio $d/m$.

All choices from Section~\ref{sec:elliptic}, including
$d,m,q,h,u_0$ and $\rho,\rho_*,\Delta,x_i,x_0,\eta$, remain
fixed. The collision takes place with $\tau$ fixed, before the
matrix limit in $\tau$.

\subsection{Collision on a fixed surface}

The collision below belongs to the broader study of limits of fat-point
conditions; compare Evain~\cite[Sections~2 and~6]{Evain2007}. We prove
the particular local constraint directly.

For each fixed $\tau$, use the coordinates and the frame in
\eqref{eq:local-chart}. If $F\in\W_\tau$ has coefficients $g_j$,
write its scalar expression in this chart as
\[
 H_F(x,y)=\sum_{j=0}^{u_0}g_j(z_0e^x)e^{jy}.
\]
Define the finite index set
\[
 \mathcal I=\{(\ell,b)\in\Z_{\geq0}^2:
                   0\leq b<m,\quad 0\leq\ell<q(m-b)\}
\]
and the linear map $J_\tau:\W_\tau\longrightarrow\C^{\mathcal I}$
whose coordinates are
\begin{equation}
 (J_\tau F)_{\ell,b}
   =\left.\partial_x^\ell\partial_y^b H_F(x,y)\right|_{(0,0)},
 \qquad 0\leq b<m,\quad 0\leq\ell<q(m-b).
 \label{eq:jet-rows}
\end{equation}
The value bundle is expressed in one fixed local frame when these
derivatives are taken.

\begin{lemma}[Collision of the further points]
\label{lem:collision}
If $\mathsf U_r(d,m)$ holds, then $\ker J_\tau\neq\{0\}$ for
every fixed $\tau\in(0,1)$.
\end{lemma}

\begin{proof}
Fix $\tau$. Choose a sequence of ordered tuples
$(\xi_{1,\nu},\ldots,\xi_{q,\nu})$ of distinct points in the
$x$-coordinate disk, all tending to zero. By
Proposition~\ref{prop:necessary-W}, there is a nonzero
$F_\nu\in\W_\tau$ with multiplicity at least $m$ at every
$(\xi_{a,\nu},0)$.

Fix a basis of the finite-dimensional space $\W_\tau$ and normalize
the coefficient vector of each $F_\nu$ to have Euclidean norm one.
A subsequence converges to a unit vector, defining a nonzero
$F_0\in\W_\tau$: the coefficient blocks in
\eqref{eq:coefficient-spaces} form a direct sum. As the basis and chart
are fixed, coefficient convergence gives locally uniform convergence
of $H_{F_\nu}$ and all of its derivatives.

For $0\leq b<m$, set
\[
 h_{b,\nu}(x)=\partial_y^b H_{F_\nu}(x,0),
 \qquad h_{b,0}(x)=\partial_y^b H_{F_0}(x,0).
\]
Multiplicity at least $m$ at $(\xi_{a,\nu},0)$ implies
\[
 \ord_{\xi_{a,\nu}}h_{b,\nu}\geq m-b
 \qquad (1\leq a\leq q).
\]
If $h_{b,0}$ is identically zero, all of its required jets already
vanish. Otherwise choose a small closed disk centered at zero,
contained in the coordinate disk, on which $h_{b,0}$ has no zeros
except possibly at zero and has none on the boundary. For all
sufficiently large $\nu$, every $\xi_{a,\nu}$ is inside this disk,
and uniform convergence on its boundary implies that $h_{b,\nu}$
and $h_{b,0}$ have the same number of zeros there, counted with
multiplicity. The argument principle, or Rouch\'e's Theorem,
therefore gives
\[
 \ord_0 h_{b,0}\geq q(m-b).
\]
This holds for every $b<m$, and is exactly $J_\tau F_0=0$.
\end{proof}

The rows in \eqref{eq:jet-rows} describe a local ideal. If the distinct centers are
$(\xi_a,0)$ and $Q(x)=\prod_{a=1}^q(x-\xi_a)$, write a local
holomorphic function as $\sum_{b\geq0}h_b(x)y^b$. The multiplicity
conditions require $Q^{m-b}$ to divide $h_b$ for $b<m$.
When the centers collide, the corresponding limiting conditions
are divisibility by $x^{q(m-b)}$, or membership in
$(x^q,y)^m$. Their number is
\[
 |\mathcal I|=\sum_{b=0}^{m-1}q(m-b)=\frac{qm(m+1)}2.
\]
Lemma~\ref{lem:collision} establishes these conditions for a nonzero limit.
It remains to compute the matrix of $J_\tau$ in bases suited to
$\tau\to0^+$.

\subsection{A common chart and normalized theta sections}

Although the tori vary, one polydisk in $(x,y)$ is a valid chart for
all sufficiently small positive $\tau$.
Fix $R$ with $0<R<\pi/2$, and let
\[
 \delta_P=\min_{1\leq i\leq9}
           \{x_0-x_i,\,1-(x_0-x_i)\}>0.
\]
For small enough $\tau$, require
$|\log\tau|>2R$ and $\delta_P|\log\tau|>R$.
If $z_0e^x$ and $z_0e^{x'}$, with $|x|,|x'|<R$, represent the
same point of $X_\tau$, then
\[
 x'-x=a\log\tau+2\pi\mathrm{i}b\qquad(a,b\in\Z).
\]
The bounds on the real and imaginary parts force $a=b=0$.
Thus this disk maps injectively to $X_\tau$. If a point of the
disk represented a point of $P$, its real coordinate would satisfy
\[
 \operatorname{Re}x=(x_i-x_0+a)\log\tau
 \qquad\text{for some }i\text{ and }a\in\Z,
\]
whose absolute value is at least $\delta_P|\log\tau|>R$.
Hence the disk avoids $P$. The covering-space frame trivializes
$A$ there, and $v=e^y$ is injective for $|y|<R$.
Thus \eqref{eq:local-chart} gives the common polydisk $|x|,|y|<R$
around the marked point on $\Tot(A)$. We continue to use the
covering-space frames of Section~\ref{sec:elliptic}.

\begin{lemma}[Normalized theta sections]
\label{lem:theta-limit}
Fix an integer $n>0$, a real number $\sigma$, and a complex number
$\epsilon$ with $|\epsilon|=1$. Put $U=\sigma-x_0n$ and assume
$U\notin\Z$. For each integer $K$ with $U<K<U+n$, there is a
section $s_{K,\tau}$ of $D(n,\epsilon\tau^\sigma)$ obtained from
the Laurent recurrence \eqref{eq:laurent-recurrence} by setting
the coefficient of $z^K$ equal to $\tau^{-x_0K}$ and all other
residue classes modulo $n$ equal to zero. These $n$ sections form
a basis. In the local covering-space frame their expressions are
\begin{equation}
 s_{K,\tau}(z_0e^x)
  =\sum_{p\in\Z}\epsilon^{-p}
     \tau^{p(K-U)+np(p-1)/2}e^{(K+np)x}
  \longrightarrow e^{Kx}.
 \label{eq:normalized-theta}
\end{equation}
The convergence is uniform on each compact subset of the complex
$x$-plane, and the same is true after any fixed number of
$x$-derivatives. Moreover,
\[
 \Pi(z_0e^x)\longrightarrow1
\]
with the same uniform holomorphic and derivative convergence.
\end{lemma}

\begin{proof}
Because $U\notin\Z$, the interval $(U,U+n)$ contains $n$ integers,
one in each residue class modulo $n$. The chosen starting
coefficients therefore give a basis by
Lemma~\ref{lem:torus-sections}. The recurrence yields
\[
 c_{K+np}=\tau^{-x_0K}\epsilon^{-p}
             \tau^{-\sigma p+pK+np(p-1)/2}.
\]
Substitution of $z=z_0e^x=\tau^{x_0}e^x$ gives the series in
\eqref{eq:normalized-theta}.

Write $a=K-U\in(0,n)$ and $\delta=\min\{a,n-a\}>0$.
For each integer $t\geq1$, the exponent of $\tau$ in the term
with $p=t$ is
\[
 ta+\frac{nt(t-1)}2,
\]
and in the term with $p=-t$ it is
\[
 t(n-a)+\frac{nt(t-1)}2.
\]
Both are at least $\delta t$. Their quadratic growth gives normal
convergence on compact $x$-disks for each fixed $\tau$, permitting
termwise differentiation. On $|x|\leq R'$, the contribution
of all terms with $p\neq0$ to the derivative of order $c\geq0$
has absolute value at most
\[
 2e^{|K|R'}\sum_{t\geq1}(|K|+nt)^c
              \bigl(e^{nR'}\tau^\delta\bigr)^t.
\]
For sufficiently small $\tau$, the quantity in parentheses is
at most $1/2$. The sum is then bounded by a constant, depending
only on $R',c,n,K$, times $\tau^\delta$, and tends to zero.
The $p=0$ term is $e^{Kx}$. This proves the asserted convergence
of the sections and all fixed derivatives on any compact disk.

For the product, put $a_i=x_0-x_i\in(0,1)$. The factors of
$\Theta(\tau^{a_i}e^x)$ differ from $1$ by terms whose total
absolute value on $|x|\leq R'$ is at most
\[
 e^{R'}\left(\sum_{p=0}^\infty\tau^{p+a_i}
                  +\sum_{p=1}^\infty\tau^{p-a_i}\right)
 =\frac{e^{R'}}{1-\tau}
                  \bigl(\tau^{a_i}+\tau^{1-a_i}\bigr).
\]
Summing over the nine indices gives a quantity $S_{R'}(\tau)$
tending to zero. The elementary product bound
\[
 \left|\prod_\nu(1+u_\nu)-1\right|
       \leq\exp\left(\sum_\nu|u_\nu|\right)-1
\]
holds for an absolutely summable family by passage from finite
products. It bounds $|\Pi(z_0e^x)-1|$ uniformly by
$\exp(S_{R'}(\tau))-1$. The product converges normally on
compact disks for each $\tau$, so its expressions are
holomorphic. Applying the uniform bound on a slightly larger
disk and then using Cauchy's derivative estimates proves
convergence of every fixed derivative on the original disk.
\end{proof}

\subsection{Bases for the polynomial spaces}

We now apply Lemma~\ref{lem:theta-limit} to the actual coefficient
spaces of $\W_\tau$. For $0\leq j\leq m$, their bundle is
\[
 T\otimes A^{m-j}
  =D\bigl(h,(-1)^{m-j}
      \tau^{(d-3m)\eta+(m-j)\rho_*}\bigr).
\]
Since $\eta-3x_0=(\rho_*-\Delta)/3$, its interval parameter is
\[
 U_j=\frac h9(\rho_*-\Delta)+(m-j)\rho_*.
\]
For $m<j\leq u_0$ with $d-3j>0$, the space in
\eqref{eq:coefficient-spaces} is the image of the sections of
\[
 L^{d-3j}=D\bigl(n_j,\tau^{(d-3j)\eta}\bigr),
 \qquad n_j=3(d-3j)=h-9(j-m)>0,
\]
under multiplication by $\Pi^{j-m}$. The interval parameter for
this source bundle is
\[
 U_j=\frac{n_j}{9}(\rho_*-\Delta).
\]
These parameters are irrational. Indeed, after substituting
$\rho_*=3(\sqrt r-3)$, the coefficient of $\sqrt r$ in either
expression is the positive integer $d-3j$, and the remaining
terms are rational because $\Delta\in\mathbb Q$. In particular,
\begin{equation}
 U_j\notin\Z\qquad
 (0\leq j\leq u_0,\ d-3j>0).
 \label{eq:nonintegral}
\end{equation}

The finite exponent set is
\begin{equation}
\begin{split}
 \mathcal S={}&
  \bigcup_{j=0}^{m}
    \left\{(j,K):K\in\Z,\quad
        U_j<K<U_j+h\right\}
 \\
 &{}\cup
  \bigcup_{\substack{m<j\leq u_0\\d-3j>0}}
    \left\{(j,K):K\in\Z,\quad
        U_j<K<U_j+n_j\right\}
 \\
 &{}\cup
  \left\{(j,0):m<j\leq u_0,\quad d-3j=0\right\},
 \qquad n_j=h-9(j-m).
\end{split}
\label{eq:exponent-set}
\end{equation}
All indices $j$ in this definition are integers. The last set is
empty unless $3$ divides $d$, and otherwise consists of the single
pair $(d/3,0)$. Because $d>3m$, that pair has $j>m$.

For a pair $(j,K)$ in the first union, choose the normalized
section from Lemma~\ref{lem:theta-limit} in
$H^0(B,T\otimes A^{m-j})$ as the coefficient of $v^j$.
For a pair in the second union, choose the normalized section in
$H^0(B,L^{d-3j})$, multiply it by $\Pi^{j-m}$, and use the result
as the coefficient of $v^j$. For the last union, use
$\Pi^{j-m}$ as the coefficient of $v^j$: the source space is
$H^0(B,L^0)=\C$, with basis the constant section $1$.
Denote the resulting fiber polynomials by $E_{j,K,\tau}$.

Each choice gives a basis of its coefficient block. For $j>m$,
multiplication by $\Pi^{j-m}$ is injective, since $\Pi$ is nonzero
on $B\setminus P$, and its image is the block by definition.
The different powers $v^j$ form a direct sum. Hence
\[
 \{E_{j,K,\tau}:(j,K)\in\mathcal S\}
\]
is a basis of $\W_\tau$ for every $\tau\in(0,1)$. There are no
duplicate pairs, and the index set is independent of $\tau$.

The block dimensions also give the useful identity
\[
 |\mathcal S|=\dim\W_\tau
 =\binom{d+2}{2}-9\binom{m+1}{2}.
\]
To see this, sum the dimensions $3(d-3j)$ of
$H^0(B,L^{d-3j})$ for $d-3j>0$ and add one when $3$ divides $d$.
The result is
$\binom{d+2}{2}$. For $j<m$, the coefficient condition in
\eqref{eq:nine-zeros} reduces the block dimension by $9(m-j)$;
their sum is $9\binom{m+1}{2}$.

In the common local chart, the scalar expressions of this basis
satisfy
\[
 H_{E_{j,K,\tau}}(x,y)\longrightarrow e^{Kx+jy}
 \qquad((j,K)\in\mathcal S)
\]
locally uniformly, together with all fixed mixed derivatives.
For the first union this is \eqref{eq:normalized-theta} times
$e^{jy}$. For the second union, the additional factor
$\Pi(z_0e^x)^{j-m}$ tends to $1$ with all derivatives by
Lemma~\ref{lem:theta-limit}; the exponent $j-m$ is fixed.
The same product assertion proves the claim for the degree-zero
case with $K=0$. As $\mathcal S$ is finite, all of these
convergences can be taken on one fixed smaller polydisk and for
the same sufficiently small values of $\tau$.

\subsection{The limiting jet matrix}

Evaluation of monomials in exponent pairs is a useful way to express
interpolation rank. Dumnicki~\cite[proof of Proposition~13]{Dumnicki2007} uses
this principle to relate a one-point multiplicity matrix to polynomial
evaluation on a diagram. Here the normalized theta sections produce
the exponent pairs directly, with the derivative indices prescribed
by the preceding collision.

\begin{proposition}
\label{prop:matrix-limit}
Choose fixed orderings of $\mathcal I$ and $\mathcal S$, and
represent $J_\tau$ in the basis $E_{j,K,\tau}$ by the matrix
$B_\tau$. These matrices have one fixed size and converge
entrywise as $\tau\to0^+$ to the matrix
\begin{equation}
 B_0=\bigl(K^\ell j^b\bigr)_
              {\substack{(\ell,b)\in\mathcal I\\(j,K)\in\mathcal S}}.
 \label{eq:limiting-matrix}
\end{equation}
If $B_0$ has full column rank, then $J_\tau$ is injective for all
sufficiently small positive $\tau$.
\end{proposition}

\begin{proof}
The basis construction gives $|\mathcal S|$ columns for every
$\tau$, and the row set is the fixed set $\mathcal I$.
The mixed-derivative convergence just established gives
\[
 (B_\tau)_{(\ell,b),(j,K)}
  =\left.\partial_x^\ell\partial_y^b
          H_{E_{j,K,\tau}}(x,y)\right|_{(0,0)}
  \longrightarrow K^\ell j^b.
\]
Zeroth derivatives contribute the factor $1$, also when $K=0$
or $j=0$. This is \eqref{eq:limiting-matrix}.
If $B_0$ has full column rank, some minor of size
$|\mathcal S|$ has nonzero determinant. The corresponding
determinant of $B_\tau$ converges to it and is therefore nonzero
for all sufficiently small positive $\tau$. Thus $B_\tau$ has
full column rank, and its linear map $J_\tau$ is injective.
\end{proof}

The factors $\tau^{-x_0K}$ may tend to zero or infinity, but for each
$\tau>0$ they give invertible basis changes and preserve rank.
Lemma~\ref{lem:collision} supplies a nonzero kernel separately for
each fixed $\tau$; Proposition~\ref{prop:matrix-limit} compares
numerical matrices in the explicit bases. Thus no continuous choice
of kernel vectors, or algebraic family of plane embeddings across
$\tau=0$, is required. To contradict Lemma~\ref{lem:collision}, it remains
to prove full column rank of the corresponding limiting matrix for a
sufficiently large common multiple of $(d,m)$.
